\section{Congruence lengths and the large-prime row}

% BEGIN G79E1
\paragraph{G79E1. The finite and extended forms of the congruence estimate.}
Fix an original datum $D$, a localization $i$ of residue characteristic
$\ell$, a point $\pi\in Z_i$, and a rank-$r$ representation $\rho$ over
the coefficient DVR $O$. Put $e=\operatorname{length}_O(O/\ell O)$.
It is positive and finite since $\ell$ is a nonzero nonunit of $O$.
Let $I_\pi$ be the ideal of values at $\pi$ of Hecke functions supported
on $Z_i$. Write $c=\operatorname{length}_O(O/I_\pi)$,
$a=\operatorname{length}_O\operatorname{Sel}_\pi(\rho)$, and
$h=v_\ell(h_E)$. For $0<j<t$ let $u_j$ be the original extended
Tate--Shafarevich length of the $j$th twist. Put
$z_P=\operatorname{length}_O(O/(\zeta_P-1))$ at $P\in Q_i$.
In extended lengths the native congruence inequality is exactly
\[
 c\le a+e\sum_{0<j<t}u_j+eh+t\sum_{P\in Q_i}z_P.                 \tag{G79.1}
\]
Indeed every value in $I_\pi$ is achieved because evaluation maps
constants onto $O$. A generator of $I_\pi$ realizes the infimum of
the principal-quotient lengths of the values, and every other value
generates a subideal. This identifies the native $v_{\rm CC}$ with
$c/e$, as in LN3--LN5. The Selmer term is $a/e$, the class term is
$h$, and the auxiliary-prime correction is $t\sum_P z_P/e$.
The lengths $u_j$ enter with their given normalization. Multiplying
by $e$ in the extended nonnegative reals gives (G79.1).

Suppose $\zeta_P\ne1$ for every $P\in Q_i$, and suppose the rational
inequality obtained by dividing (G79.1) by $e$ is established whenever
$a$ and every $u_j$ are finite. Then (G79.1) holds without those
finiteness assumptions. An infinite $a$ or $u_j$ makes the right side
infinite since $e>0$. In the remaining case all $z_P$ are finite.
Also $I_\pi\ne0$: for each point outside $Z_i$ choose a Hecke
function distinguishing it from $\pi$, subtract its value at that
point and take the product over the finite complement. This product
vanishes there and has nonzero value at $\pi$, also when the
complement is empty. Thus $c$ is finite. The specified finite
$u_j$ agree with the native rational row values by the datum's
Tate--Shafarevich length identity. Clearing the positive denominator
in the rational inequality proves (G79.1), including its exact integer
coefficients. This proves the original extended-from-rational theorem.
% END G79E1

% BEGIN G79E2
\paragraph{G79E2. From the extended estimate to zero additional error.}
Assume the original $\mathrm{CongBound\_at}(D,L)$,
$\mathrm{WeightParam}(D)$, $\mathrm{PtFam}(D,R,\rho)$,
$\mathrm{KillAt}(\rho)$ and $\mathrm{ShaFinOdd}(D)$.
Let $p_*=\ell_*(L)\ge3$. Take a finite set $E_L$ whose underlying
set equals $\mathrm{BaseSizes.Exc}(L)$. Then
$\mathrm{StarE}(L,\rho,E_L)$ holds.

Indeed take an arbitrary twisted $i$, with $\ell>p_*$,
$\ell\notin\{g_0,g_1\}\cup E_L$, and $\pi\in Z_i^{\rm good}$.
The good-point identity supplies both $\pi\in Z_i$ and the witness
$\pi$ for $\mathrm{HG}$ at $i$. The required archimedean
$\mathrm{WeightTie}$ follows from $\mathrm{WeightParam}$ on the
given embedding and its complex conjugate. More explicitly the two
parameter multisets have identical cloud and low terms and identical
length-$t$ strings after reindexing $c=1,\ldots,t$ by $c=j+1$.
Use shift $-1/2$ on the given embedding and negate the entire shifted
multiset on the conjugate embedding. The pointwise family gives
$\mathrm{GoodR}$ for $R_i$ and a continuous complement of $R_i(\pi)$
with precisely the prescribed $t$ character lines; the same splitting
over the fraction field is $\mathrm{IsRankRPart}$ (DFD6--DFD7).
The guard on $E_L$ is the required guard on the arithmetic exceptional
set, so $\mathrm{CongBound\_at}$ gives (G79.1) for this actual complement.

Killing gives finite $a$ and
$t\sum_P z_P/e\le1/\ell^2$. Here $\zeta_P$ is the twist value
on the chosen tame generator. At $P\in Q_i$ the base character is
unramified, so this equals the native first-character value there.
It differs from $1$: the twist ramifies at $P$, is of $\ell$-power
order and kills wild inertia since $P\nmid\ell$, so it is nontrivial
on a generator of tame inertia. Thus every $z_P$ is finite.
The prime $\ell>3$ is odd, so all the $u_j$ are finite. Dividing
(G79.1) by $e$ and applying the auxiliary-prime bound gives
\[
 v_{\rm cc}(i,\pi)\le s_\pi(\rho_i)+
 \sum_{1\le j<t}u_j+v_\ell(h_E)+1/\ell^2.
\]
This is exactly $\mathrm{Bd}(\rho,i,\pi,0)$, including the
$1/\ell^2$ term built into $\mathrm{Bd}$. Only the additional
error parameter is zero. This proves the original
\texttt{localBoundZeroAtLargePrimes\_\allowbreak of\_\allowbreak congruenceBound}.
% END G79E2

% BEGIN G79E3
\paragraph{G79E3. The unshifted jump windows at every letter.}
Take any original family and a member with its parameter choice at
$x>\mathcal F.x_1$. Use $p_*,b,b',w'',S_c,\eta,\mathcal B$ of X78A1
and put $B=b/p_*$ and $T=(t-1)/2$.
Every $W\in\mathcal B$ satisfies
\[
 2B+t/2+3\le W\le b.
\]
The lower bound follows from $W\ge w''\ge16\eta$ and
$12B+2t+8\le\eta$. For the upper bound, $b,w''\le b$.
The other centers are at most $9b'/4$, as in X78A1. The denominator
defining $p_*$ is at least $t>3m$, so $3p_*<b$ and
$9b'/4<3p_*<b$.

Fix $1\le j\le m$, a letter with embedding $\tau$, and its orientation.
Let $F_l,F_c$ be respectively the integral points of its native low
and cloud sets. They are finite as these sets are images of finite
lists or their reflections in $T$. Reflect a downward letter by
$z\mapsto t-1-z$, and leave an upward letter fixed; put a prime on
the resulting points and sets. In either orientation one has
\[
 2B+t/2+3\le |z'-T|\le\eta\quad(z'\in F_l'),\qquad
 W\le z'-T\le W+S_c+\eta+1\quad(z'\in F_c')
                                                               \tag{G79.2}
\]
for one center $W$ depending on the letter. In fact
$z'-T=e_{j-1}(\tau)-b_d(\tau)$ on the low list, so its two
bounds are the size condition on the low differences and their
separation in the parameter choice. On the cloud list it is
$e_{j-1}(\tau)+v_d(\tau)$ with $0\le v_d\le S_c$.
The top, thick and thin size conditions place $e_{j-1}(\tau)$ in
$[W,W+\eta+1]$ for $W=b_\tau+\lfloor b'/x\rfloor$, $b$ or
$w''$ respectively. The extra one covers the floor in the top case.
All statements concern the integral points even when a defining
list had rational entries or repeated entries.
% END G79E3

% BEGIN G79E4
\paragraph{G79E4. Shifts, carries and residue avoidance for a cyclic list.}
Fix any prime $\ell>p_*$ outside the exceptional set of X78A2, any
positive cycle length $n$ and any choices of embeddings and
orientations around $\mathbb Z/n$. At each letter work first with
the upward coordinates of G79E3. Include the endpoints of the
character string by setting
\[
 m_0=\min(F_l'\cup\{0\}),\quad M_0=\max(F_l'\cup\{t-1\}),\quad
 M=\max(F_c'\cup\{M_0\}).
\]
Set $w_+=1-m_0$, $d=M_0-m_0+1$, $h=M-m_0+2$.
Let $U_+$ contain each $z'+w_+$, $z'\in F_l'\cup F_c'$, once.
For upward orientation take $w=w_+$, $U=U_+$. For downward
orientation take $w=h-w_+-(t-1)$ and
$U=\{h-y:y\in U_+\}$. In either case the resulting multiset is
the translate by $w$ of the original union, each integer counted once.

Since $t\le\eta$, the low endpoints satisfy $T-\eta\le m_0\le 0$
and $t-1\le M_0\le T+\eta$. The cloud points lie above $t-1$
by (G79.2). Hence $1\le d\le2\eta+1$, all $U_+$ lie in $[1,h-1]$,
and $[w_+,w_++t-1]\subset[1,d]\subset[1,h-1]$. Reflection places
the full downward character string in $[h-d,h-1]$ and all its jumps
in $[1,h-1]$. In particular both prescribed shifts are positive.
For either orientation (G79.2) and $W\le b$ give
$h\le H=b+S_c+2\eta+3$, uniformly over the cycle.

At letter $i$ its integer carry is
$\sigma_i=\lfloor\sum_{k\ge1}h_{i+k}/\ell^k\rfloor$, with indices
read cyclically. The series is dominated by the geometric series
$H/(\ell-1)$. The two size inequalities of X78A1 give
$S_c<p_*$ and $2\eta+2B+3<p_*$, hence
$S_c+2\eta+2B+3<2p_*<b$. It follows that
$H\le 2B(p_*-1)$ and $0\le\sigma_i\le s=\lfloor2B\rfloor$.
In particular $d+4\sigma_i+6\le2\eta+8B+7<\ell$.
For any original upward-coordinate jump $z'$ and
$v\in\{0,\ldots,t-1\}$, (G79.2) yields
$|z'-v|\ge |z'-T|-T\ge 2B+7/2>\sigma_i+2$.
Translations and reflection preserve distances, so every jump in
the specified low window has the required separation.

Finally take an upward-coordinate $y=z'+w_+$ outside $[1,d]$.
Every low point belongs to that window, hence $z'\in F_c'$.
For any integer $a$ with $-2\sigma_i-2\le a\le d+2\sigma_i+3$,
substitution of $d,w_+$ gives
\[
 W-\eta-2s-3\le z'-M_0-2\sigma_i-3
     \le y-a\le z'-m_0+2\sigma_i+3
     \le W+S_c+2\eta+2s+4 .
\]
By X78A2, $\ell$ divides no integer in this window. Therefore
$y\not\equiv a\pmod\ell$. On a downward letter the native
residue test uses $h-y_{\rm down}$, exactly the upward coordinate
just treated. This proves every inequality and residue clause of
$\mathrm{PositionHyps}$ with window width $d$.
% END G79E4

% BEGIN G79E5
\paragraph{G79E5. The original uniform local position-data theorem.}
Use DFD14 at $C=1$ to choose one family cutoff beyond which $q,d$
and every prime in the host's small set are at most $p_*$.
For each member at that cutoff with $x>\mathcal F.x_1$, any prime
$\ell>p_*$ consequently avoids these primes and satisfies
$\ell\ge5$ and $t+1<\ell$ by X78A1. These guards hold before
excluding any further primes. Now assume $\ell$ avoids the arithmetic
exceptional set. For every $1\le j\le m$, every $n\ge1$ and all
orientation and embedding functions on $\mathbb Z/n$, use G79E4
with the window type equal to the prescribed orientation.
For every integer $a$ and every letter $i$,
\[
  \operatorname{count}_{U_i}(a+w_i)
  =\begin{cases}
     1,&a\in\mathrm{JW}(D,\mathrm{up}_i,j,\tau_i),\\
     0,&\text{otherwise}.
   \end{cases}
\]
Indeed $\mathrm{JW}$ is the union of the native low and cloud
sets, and G79E4 translated their integral union as a set.
Together with $\mathrm{PositionHyps}$ this proves
\texttt{exists\_localPositionData\_}\allowbreak\texttt{off\_exceptional\_primes},
including its quantifiers under the stated $\mathrm{JumpConditions}$
premise. The construction itself uses only the size and parameter
conditions of the family. Identifying the jumps of a crystalline
representation with these integer sets is a further arithmetic assertion.
% END G79E5

% BEGIN G79E6
\paragraph{G79E6. The class number and cyclotomic contributions.}
Continue with the native DVR $O$, fraction field $K$ and $e$ of
G79E1. The $t$ character blocks in the count have ratios
$\chi_{j+1}/\chi_1=\chi_{\rm cyc}^{-j}$, $0\le j<t$.
For $j=0$ consider an $O$-module of continuous cocycles into $K/O$
with zero restriction to inertia at every finite place. The action
is trivial, so these are additive homomorphisms. Under global Artin
reciprocity each factors through the ideal class group of $E$:
the finite local units map onto the inertia images, the archimedean
components of this CM field are connected, and quotienting the
ideles by these factors and principal ideles gives the ideal class
group. The maps have finite image by compactness and discreteness.
A cyclic factor of order $m$ in the class group contributes at most
the length of $(K/O)[m]$, which is $e v_\ell(m)$ by the
isomorphism $O/mO\simeq m^{-1}O/O$. Adding over the cyclic factors
gives the bound $e v_\ell(h_E)$, including when $\ell\nmid h_E$.

For $1\le j<t$ let $C$ be an $O$-submodule of continuous
$K/O(-j)$-valued cocycles. Suppose its elements vanish on inertia
away from $\ell$ and are coboundaries on the decomposition groups
above $\ell$. In the native local-coordinate description the latter
hypothesis is membership in $\mathrm{BD}$ after localization: the
decomposition group and local Galois group are identified by the
bridge, whose coefficient-field isomorphism carries the lattices
onto each other. Pull a local cobounding vector back under this
identification. Then
\[
 \operatorname{length}_O C\le
 \operatorname{length}_O(C\cap B^1)+e\,u_j .
\]
Choose a $\mathbb Z_\ell$-basis of $O$, of rank $d_0$.
It extends to a basis of $K$ and identifies $K/O$ with
$(\mathbb Q_\ell/\mathbb Z_\ell)^{d_0}$. The twist is a scalar
in $\mathbb Z_\ell^\times$, so every coordinate is a cocycle for
the native $-j$ action. Each descends to the quotient unramified
away from $\ell$: the action there is trivial on the closed normal
subgroup generated by the other inertias, and a continuous cocycle
zero on those inertias vanishes on this subgroup by the cocycle law.
Above $\ell$ each coordinate is locally a coboundary, giving
$d_0$ classes in the original Tate--Shafarevich group. Conjugating
the chosen decomposition group changes a cobounding vector $m$ to
$g m-z(g)$ for the translated group, so this is independent of the
place lifts. If all the classes vanish, pick one global cobounding
vector per coordinate and assemble them in $K/O$. This proves the
injection, after restriction of scalars, of $C/(C\cap B^1)$ into
$d_0$ copies of the native Tate--Shafarevich group.
If $u_j$ is finite the target has $\mathbb Z_\ell$-length $d_0u_j$.
Each $O$-composition factor has $\mathbb Z_\ell$-length
$f_0=[O/\mathfrak m:\mathbb F_\ell]$, and $d_0=e f_0$ by
computing the length of $O/\ell O$ in both rings. Dividing by
$f_0$ proves the inequality. Infinite $u_j$ gives it immediately.
% END G79E6

% BEGIN G79E7
\paragraph{G79E7. The unit root difference eliminates global coboundaries.}
Suppose a native orthogonal decomposition with its integral root
data is given, with chosen lift $\pi'$ of $\pi$. In its finite
matrix family write $A$ for the algebra spanned over its Hecke
algebra $T'$ by the matrices $R(g)$. Write $E_1$ for the integral
root polynomial cutting out the first character and $\delta$ for the
product of its differences from the other roots. Thus $E_1\in A$,
$\delta\in T'$, $E_1^2=\delta E_1$, and at $\pi'$ the blocks of
$E_1$ are $\delta(\pi')$ on $\chi_1$ and zero elsewhere.
These identities belong to the orthogonal decomposition, with
$E_1\in A$ also following by multiplying its integral linear factors.
Assume $\delta(\pi')$ is a unit. Then the native rank and character
cocycle modules contain no nonzero global coboundary.

Indeed put $\alpha=\chi_1$ and
$y_g=(\alpha(g)^{-1}R(g)-1)E_1$. A generator of a native cocycle
module has the form $z(g)=\phi(y_g)$, where $\phi$ is an $O$-linear
map on the finite matrix family, intertwining the left $A$-action
on $Y=\sum A y_g$ with the selected block $\psi$ at $\pi'$.
The defining support condition on $\phi$ permits evaluation on any
specified subset of points. These generators and their $O$-linear
combinations are continuous cocycles: the identity
$y_{gh}=y_g+\alpha(g)^{-1}R(g)y_h$ gives the cocycle law for
$M(g)=\alpha(g)^{-1}\psi(R(g))$. Continuity follows from that of
the finite matrix family and of $\alpha$, since $\phi$ on the finite
free source has discrete torsion target and kills a power of the
uniformizer times the source. The block at $\pi'$ is the indicated
$\rho$ or $\chi_{j+1}$ by applying the decomposition to their standard
coordinate vectors.

Write $E_1=\sum f_j R(g_j)$ with $f_j\in T'$ and set
$b_j=\alpha(g_j)f_j$. Then
$Y_0=\sum b_j y_{g_j}=(\delta-\sum b_j)E_1$ belongs to $Y$ and
$E_1Y_0=\delta Y_0$. For any block on which $E_1$ acts as zero,
intertwining these two sides shows
$\delta(\pi')\sum b_j(\pi') z(g_j)=0$. This conclusion is linear
in $z$. If $z(g)=(M(g)-1)v$ is a global coboundary in that block,
then $\sum b_j(\pi')M(g_j)=\psi(E_1)=0$, while
$\sum b_j(\pi')=\delta(\pi')$ by reading the same expansion
on the first character. The relation therefore gives
$\delta(\pi')^2 v=0$, hence $v=0$. On the first character
$M(g)=1$ and all coboundaries are already zero. This proves the
zero-length coboundary assertions in the unit case used in the global count.
% END G79E7

% BEGIN G79E8
\paragraph{G79E8. Tame restrictions of the character blocks.}
Use the character cocycles from G79E7 and a place
$P\in S_i\cup Q_i$, hence $P\nmid\ell$. Each ratio character on
inertia is trivial, since its cyclotomic factor is unramified
there. Each restricted cocycle is therefore an additive homomorphism.
It kills wild inertia: its image is a finite $\ell$-group by
continuity and compactness, and wild inertia is pro-residue-characteristic.
It is determined by its value on the fixed tame generator $\gamma_0$.

Under the native quadratic inertia relation
$(R(\gamma)-1)(R(\gamma)-\alpha(\gamma))=0$ at every selected
point, $(R(\gamma)-1)y_\gamma=0$ there. Applying the defining
support and intertwining conditions of $\phi$ gives
$(\chi_{j+1}(\gamma)-1) z(\gamma)=0$. In particular
$z(\gamma_0)\in(K/O)[\zeta_P-1]$, since every character agrees
with $\alpha$ on inertia. Evaluation is injective on the restriction
module. Thus its length is at most $z_P$ when $\zeta_P\ne1$ by
the principal-quotient calculation in G79E6. It vanishes when
$\zeta_P-1$ is a unit. Under the alternative native polynomial
condition at $P\in S_i$, let $P_0$ be its polynomial over the Hecke
algebra. The matrix $A_0=P_0(R(\gamma_0))$ belongs to $A$ by
expanding its powers with the given Hecke coefficients, and the
relation says $A_0 y_\gamma=0$ on the selected points for every
inertia element. Its block scalar is $P_{0,\pi}(\zeta_P)$,
a unit by hypothesis. Applying $\phi$ again gives $z(\gamma)=0$.
All three conclusions pass to spans by linearity.

Away from $S_i\cup Q_i$ and $\ell$ the native unramified-family
condition says $\alpha(\gamma)^{-1}R(\gamma)=1$ on the selected points
for inertia elements, so $y_\gamma=0$ there. Hence all the corresponding restrictions, including
those on the rank block, vanish there. At $Q_i$ impose the quadratic
condition and $\zeta_P\ne1$ for every $P$. For any character block
$C_j$, let $C'_j$ be the kernel of the joint inertia restriction
at these places. Since $C_j/C'_j$ embeds in the direct sum of
their individual restriction images,
\[
 \operatorname{length}_O C_j\le
 \operatorname{length}_O C'_j+\sum_{P\in Q_i}z_P.
\]
If every $P\in S_i$ satisfies either the quadratic relation with
$\zeta_P-1$ a unit or the stated polynomial condition, the cocycles
in $C'_j$ vanish on all inertia away from $\ell$.
% END G79E8

% BEGIN G79E9
\paragraph{G79E9. The remaining determinant and local inputs in the rational count.}
Fix the original $i,\pi,R,\rho$ of the congruence problem.
Use the full set of points, a unit-root-difference decomposition,
and the unramified, quadratic and polynomial hypotheses of G79E8.
Assume global Artin reciprocity and the following two estimates or
inclusions on its original cocycle modules. First, the determinant
bound for this decomposition gives
\[
 c\le\operatorname{length}_O C_\rho+
             \sum_{0\le j<t}\operatorname{length}_O C_j .
\]
Second, the localizations of those cocycles at every $w\mid\ell$
belong respectively to the Bloch--Kato finite part for $C_\rho$,
to the unramified cocycles for $C_0$, and to coboundaries for
$C_j$, $j>0$. These hypotheses imply (G79.1).

Indeed the classes of $C_\rho$ belong to the original Selmer
group: above $\ell$ this is the stated finite condition, away from
$S_i,Q_i,\ell$ the inertia restrictions vanish by G79E8, and
the remaining allowed places impose no further condition.
The kernel of passage to classes is zero by G79E7, so
$\operatorname{length}_O C_\rho\le a$. On the first-character
block, being unramified above $\ell$ implies actual vanishing on
inertia there, because the ratio action is trivial. Hence G79E6
bounds $C'_0$ by $eh$. On the other character blocks G79E6 bounds
$C'_j$ by $e u_j$ since G79E7 kills the coboundary term.
G79E8 adds the same $\sum_P z_P$ to each of the $t$ blocks.
Sum and apply the indicated determinant estimate. For the full
set of points the evaluated ideal in that estimate is $I_\pi$,
since both definitions evaluate Hecke functions vanishing outside
$Z_i$. This yields precisely (G79.1). G79E1 supplies the original
extended formulation, including its infinite cases.

The formal route supplies the second input by localizing the global
cocycles in the local cut spans. Its rank-block condition says that
the quotient of that span by its intersection with the finite part
has length zero. Thus the quotient itself vanishes.
The character-block conditions refer to those same cut spans.
G79E5 supplies their numerical position data; producing their
local Galois and cohomological conditions is a separate step.
% END G79E9

% BEGIN G79E10
\paragraph{G79E10. The precise large-prime field on the chosen family hub.}
Under J78H1 let a given member also satisfy
$\mathrm{CongBound\_at}(D,L)$. On its very same hub one obtains
$\mathrm{Concl.large}$. Indeed the member's core facts give
$\mathrm{WeightParam}$, the canonical setup gives
$\mathrm{PtFam}(D,X.R,X.\rho)$, and setup and the family premise
give killing and odd finiteness. The small-prime estimate gives
$p_*\ge3$; X78A3 identifies $X.\mathrm{exc}$ with the exact
exceptional set. Apply G79E2. A uniform proof of
$\mathrm{CongBound\_at}$ for every sufficiently large member
satisfying its native facts therefore also supplies the large-prime
field on that same hub at the maximum of its cutoff and that of J78H1.

Here the required antecedent quantifies over every localization with
$\ell>p_*$ avoiding both $g$'s and the arithmetic exceptional set,
with the prescribed weight tie and $\mathrm{HG}$, every native
good representation family, every $\pi\in Z_i$ and every continuous
rank-$r$ complement. Its quantifiers include points beyond the
twisted good points on which $\mathrm{StarE}$ is evaluated.
The original construction uses the member's jump and image
conditions, the integral orthogonal decomposition with a unit
root difference, the determinant bound and the local cut-cocycle
conditions of G79E9. DFD13 supplies the eventual parameter
inequalities, L78W13 supplies the Wach comparison, and G79E5
settles the numerical position tables. G79E11--G79E13 below deduce the determinant count from the given
unit-root-difference decomposition and the off-support hypothesis.
Constructing this decomposition and proving the required local and
image conditions from the member's facts are further inputs.
% END G79E10

% BEGIN G79E11
\paragraph{G79E11. Determinant coefficients descend to the original evaluation ideal.}
Take the orthogonal decomposition of G79E7 on the full set of points,
and assume the native off-support condition: for $x\notin Z_i$,
$\alpha^{-1}R_x$ has no fixed vector over $K$. Use the algebra $A$,
the root cover $p:X'\to X_i$, its Hecke algebra $T'$ and its
distinguished point $\pi'$ from that decomposition. The submodule
$Y=\sum_g A y_g$ of the finite matrix family is Noetherian over $O$,
so it is generated over $A$ by finitely many $y_{g_q}$, $1\le q\le s$.
Write
\[
 \mathcal R=\{(a_q)\in A^s:\textstyle\sum_q a_q y_{g_q}=0\}.
\]
Take any finite list of square $s$-by-$s$ matrices $B_k$ over $A$
whose rows belong to $\mathcal R$, and use independent commuting
variables $u_k$. At a point of the cover consider the determinant
of the ordinary matrix obtained by replacing each entry of
$B=\sum_k u_k B_k$ with its matrix of size $n_V$. All coefficients
of this determinant are in $T'$, and their values at $\pi'$ lie
in $I_\pi$.

Here is a proof of coefficient membership using the determinant
condition of the decomposition. That condition places the coefficients
of $\det(\sum v_i a_i)$ in $T'$ for any finite $a_i\in A$.
After polynomial substitution it applies to a single $n_V$-square block
with entries in the algebra $A$ extended by scalar indeterminates;
it also applies to each coefficient in an auxiliary formal variable
by truncating the series to the desired order. Form $1+zB$ and
eliminate its first block column. Its upper left block has constant
term $1$, so its inverse has coefficients in $A[u_k][[z]]$ by
the geometric series. The Schur complement again has constant term
the identity and coefficients in this algebra. At every point of
$X'$ the determinant is the product of the pivot determinant and
the determinant of the Schur complement, by block triangular
elimination. Repeat for all blocks. The single-block assertion shows
that every coefficient of the product lies in $T'$. Taking the
coefficient of $z^{s n_V}$ gives the desired determinant. This
argument uses no division by an integer.

For $p(x')=x\notin Z_i$ the rank-one operator $E_1(x')$ has a
nonzero image vector. Some $\alpha(g)^{-1}R_x(g)$ moves that
vector by the off-support condition, so some $y_g(x')v\ne0$.
For that $v$ at least one $y_{g_q}(x')v$ is nonzero since the
$y_{g_q}$ generate $Y$. The column consisting of these $s$ vectors
is therefore a nonzero common kernel vector of all the $B_k(x')$.
The determinant polynomial vanishes at $x'$. Consequently each of
its coefficients $f\in T'$ vanishes on every fiber above $X_i\setminus Z_i$.

The orthogonal decomposition gives a basis $(b_a)$ of $T'$ over
the original Hecke algebra $T_i$. Write $f=\sum (p^*f_a)b_a$,
$f_a\in T_i$. Over $K$, the original algebra is $K^{X_i}$,
by the same point-separating products as in G79E1, divided by
their nonzero value at the chosen point. After extending scalars
$T'_K$ still embeds in $K^{X'}$, and multiplication by the
idempotent supported at $x$ projects onto its fiber. Thus the
restrictions of the $b_a$ are linearly independent on each fiber:
a relation there would give a $T_{i,K}$-relation supported at $x$.
Hence each $f_a$ vanishes on $X_i\setminus Z_i$. Evaluating the
basis expansion gives $f(\pi')=\sum f_a(\pi)b_a(\pi')\in I_\pi$.
At $\pi'$ simultaneous decomposition into the rank block and the
$t$ character blocks factors the same determinant into the product
of their block determinants, by conjugating and reordering rows
and columns together. Every coefficient of this product is therefore
in $I_\pi$.
% END G79E11

% BEGIN G79E12
\paragraph{G79E12. The determinant bound from the block relation lattices.}
In the setting of G79E11 label the blocks by $b$, with dimensions
$d_b=r,1,\ldots,1$ and homomorphisms
$\psi_b:A\to M_{d_b}(O)$. These are algebra homomorphisms
because they are the diagonal blocks of evaluation in one fixed
basis over $K$ and preserve the integral lattices. They are
surjective by the decomposition for the rank block and by scalar
matrices for the characters. Let $N_b\subset O^{s d_b}$ be
generated by all rows of the $d_b$-by-$s d_b$ matrices
$(\psi_b(a_1)\ \cdots\ \psi_b(a_s))$, $(a_q)\in\mathcal R$.
The solution module consists of tuples in $(K/O)^{s d_b}$
annihilated by all these rows. It is
$\operatorname{Hom}_O(O^{s d_b}/N_b,K/O)$ via the dot product.

This solution module has length at most the length of the native
cocycle module $C_b$. In fact, given a solution $(v_q)$ in block
coordinates, define $\phi(\sum_q a_q y_{g_q})=
\sum_q\psi_b(a_q)v_q$ on $Y$. The relations give well-definedness
and this formula intertwines the left $A$-action. It extends
$O$-linearly to the matrix family: $(K/O)^{d_b}$ is divisible and
thus injective over the PID $O$ (on an ideal $(a)$, extend a map
by choosing an $a$-preimage of the image of $a$).
For the full set of selected points the support condition holds.
The resulting block cocycle takes value $v_q$ at $g_q$. Thus the
solution module is contained in the image of evaluation of $C_b$
at this list, giving the length inequality.

If any $N_b$ has deficient rank, its solution module contains
$K/O$ and the desired sum of cocycle lengths is infinite.
Otherwise $N_b$ has a basis of $s d_b$ rows. Smith normal form
shows its basis determinant has valuation
$u_b=\operatorname{length}_O(O^{s d_b}/N_b)$, also the length
of the solution module since $(K/O)[a]\simeq O/(a)$ for $a\ne0$.
Each chosen row of $N_b$ can be placed in any single row of the
block image of one relation, with the other rows zero. Indeed
multiply the contributing relations on the left by lifts of the
appropriate matrix units using surjectivity of $\psi_b$, multiply
by their $O$-coefficients, and add. Grouping $d_b$ such rows at a
time produces an $s$-by-$s$ matrix $B_b$ of relations whose block
image in $b$ has determinant of valuation $u_b$.

Apply G79E11 to all these matrices together. For each block $b$
let $P_b$ be its determinant polynomial in the variables of
$\sum_{b'} u_{b'}' B_{b'}$. Its minimum coefficient valuation is
at most $u_b$, by specializing its block variable to $1$ and the
others to $0$. Minimum coefficient valuations add under products:
divide each nonzero polynomial by the largest common power of a
uniformizer and reduce modulo the maximal ideal, where the product
of the resulting nonzero polynomials stays nonzero. All coefficients
of $\prod P_b$ lie in $I_\pi$ by G79E11, so
\[
 \operatorname{length}_O(O/I_\pi)\le
 \sum_b u_b\le\operatorname{length}_O C_\rho+
                 \sum_{0\le j<t}\operatorname{length}_O C_j .
\]
This gives the determinant input of G79E9 from the original
decomposition and off-support hypotheses.
% END G79E12

% BEGIN G79E13
\paragraph{G79E13. The full global count under its native local hypotheses.}
Assume global Artin reciprocity. Fix a native good family, a point
in $Z_i$, and the orthogonal decomposition for its specified
complement. Suppose there is no $\chi_1$-line off $Z_i$, the
unramified and ramified-place conditions of G79E8 hold, and above
$\ell$ the native local cut spans satisfy the zero rank-block
bound, the coboundary condition on characters indexed $2,\ldots,t$,
and the unramified condition on the first character.
Then the rational large-prime congruence bound holds whenever
$a$ and all $u_j$ are finite. This is the original
\texttt{globalCongruenceCount\_}\allowbreak\texttt{of\_rootDifference\_}\allowbreak\texttt{unit}.

To verify that the hypotheses on local cut spans apply to the
global cocycles, put $R^*(g)=\alpha(g)^{-1}R(g)$ and take one
generator $z(g)=\phi(y_g)$ of G79E7. List the columns of $E_1(x')$
for all $x'$, repeating the point as needed in the list of the
original representations at $p(x')$. Together these columns form
a vector $v$ in a finite product of copies of the representations
over $K$. The columns of $y_g$ together form $R^*(g)v-v$.
Let $Y_0$ be the $O$-span of these latter vectors. It is finitely
generated because its integral column module is finite free and
$O$ is Noetherian. It is stable since
$R^*(g)(R^*(h)v-v)=(R^*(gh)v-v)-(R^*(g)v-v)$.
Reassembling the columns and applying $\phi$ gives an
$O$-linear equivariant map on $Y_0$ with value $z(g)$ on
$R^*(g)v-v$. These are exactly the finite list, vector, stable
submodule and equivariant map in the definition of a local cut
generator. Restricting to the decomposition group and applying
the specified coefficient isomorphism gives the localization of
$z$. Taking spans therefore puts every original cocycle in the
appropriate local cut span.

In the rank block the quotient of the localized span by its
intersection with the finite part has length at most zero times
the local length of $O/(\ell)$. That local length is finite, so
the quotient vanishes and the span lies in the finite part.
The character conditions likewise pass to spans because
coboundaries and unramified cocycles are submodules. This supplies
the second input of G79E9. The orthogonal decomposition includes
the unit root difference. The off-support hypothesis gives G79E12,
which supplies the first input. Now G79E9 and division by $e$
give the stated rational inequality with the four original terms.
On an arbitrary native member the arithmetic task still needed
for $\mathrm{CongBound\_at}$ is to obtain the hypotheses of this
global count simultaneously at every guarded localization and every
allowed complement.
% END G79E13
